\documentclass[12pt]{article}                                                                                                           
\usepackage[utf8]{inputenc}                                                                                                             
\usepackage[T1]{fontenc}                                                                                                                
\usepackage{amsmath, amssymb}                                                                                                           
\usepackage{geometry}                                                                                                                   
\usepackage{xcolor}                                                                                                                     
\usepackage{lipsum}                                                                                                                     
\usepackage{titlesec}                                                                                                                   
\usepackage{fancyhdr}                                                                                                                   
\usepackage[most]{tcolorbox}                                                                                                                  
\usepackage{graphicx}                                                                                                                   
\usepackage{float}                                                                                                      
\usepackage{tikz}
\usetikzlibrary{shapes.geometric, arrows.meta}
\usetikzlibrary{arrows.meta, decorations.markings}                                                                                      
                                                                                                                                        
\geometry{a4paper, margin=2.5cm}                                                                                                        
\pagestyle{fancy}                                                                                                                       
\fancyhf{}                                                                                                                              
\rhead{Problem Set II}                                                                                                                   
\lhead{Rigid body mechanics}                                                                                                                     
\cfoot{\thepage}                                                                                                                        
                                                                                                                                        
% Fancy titles                                                                                                                          
\titleformat{\section}{\normalfont\Large\bfseries}{Exercise \thesection:}{1em}{}                                                        
\titleformat{\subsection}{\normalfont\bfseries}{Correction:}{1em}{}                                                                     
                                                                                                                                        
% Custom box for correction with breakable option
\newtcolorbox{correctionbox}{                                                                                                           
  colback=gray!5,                                                                                                                       
  colframe=black,                                                                                                                       
  fonttitle=\bfseries,                                                                                                                  
  title=Correction,                                                                                                                     
  before skip=10pt,                                                                                                                     
  after skip=10pt,
  breakable,
  enhanced
}                                                                                                                                       
                                                                                                                                        
% Fix headheight warning
\setlength{\headheight}{14.49998pt}

\begin{document}

\begin{center}
\Large\textbf{Ibn Tofail University} \\[1em]
\large\textit{Rigid body mechanics} \\[2em]
\large\textit{Problem Set II} \\[0.5em]
\end{center}

\vspace{1cm}
% ----------- EXERCISE 1 -------------
\section{}

Determine the number of degrees of freedom of the following systems:

\begin{enumerate}
    \item Two material points $M_1$ and $M_2$ moving freely on a plane while maintaining a constant distance between them.
    \item A sphere free to move in space.
    \item A sphere free to move on a plane.
    \item A sphere free to move on a line.
    \item A rod free to move on a plane.
    \item A rod with one fixed end moving on a plane.
    \item A system of two articulated rods with one fixed end.
    \item A hoop free to move on a plane.
    \item A hoop free to move on a line.
\end{enumerate}

\begin{correctionbox}
The number of degrees of freedom (DOF) is calculated as: DOF = Total parameters - Constraints.

\begin{enumerate}
    \item \textbf{Two points on a plane with constant distance:}
    
    Each point has 2 DOF in the plane (4 total). The constraint of constant distance removes 1 DOF.
    $$\text{DOF} = 4 - 1 = 3$$
    
    \item \textbf{Sphere free in space:}
    
    Position of center: 3 DOF. Rotation about center (any axis): Since a sphere is spherically symmetric, only 2 independent rotations are distinguishable (like latitude and longitude of a marked point). However, for a general rigid body sphere: 3 DOF.
    $$\text{DOF} = 3 + 3 = 6$$
    But if the sphere is perfectly homogeneous with no markings, rotations are indistinguishable: DOF = 3.
    
    \item \textbf{Sphere on a plane:}
    
    Position on plane: 2 DOF. Height fixed by contact: 0 additional DOF. Rotations: 3 DOF (roll, pitch, yaw). Constraint: contact with plane removes 1 DOF (vertical position).
    $$\text{DOF} = 2 + 3 = 5$$
    
    \item \textbf{Sphere on a line:}
    
    Position along line: 1 DOF. Rotations: 3 DOF.
    $$\text{DOF} = 1 + 3 = 4$$
    
    \item \textbf{Rod on a plane:}
    
    Position of center: 2 DOF. Orientation in plane: 1 DOF (angle).
    $$\text{DOF} = 2 + 1 = 3$$
    
    \item \textbf{Rod with fixed end on a plane:}
    
    One end is fixed: 0 DOF for position. Orientation in plane: 1 DOF (angle of rotation about fixed point).
    $$\text{DOF} = 1$$
    
    \item \textbf{Two articulated rods with one fixed end:}
    
    First rod: 1 DOF (angle from fixed point). Second rod: 1 DOF (angle relative to first rod at joint).
    $$\text{DOF} = 2$$
    
    \item \textbf{Hoop on a plane:}
    
    Position of center: 2 DOF. Orientation (rotation about perpendicular axis): 1 DOF. Rotation about its own axis: 1 DOF.
    $$\text{DOF} = 2 + 2 = 4$$
    
    \item \textbf{Hoop on a line:}
    
    Position along line: 1 DOF. Orientation (tilt): 2 DOF. Rotation about its axis: 1 DOF.
    $$\text{DOF} = 1 + 3 = 4$$
\end{enumerate}
\end{correctionbox}

\newpage

% ----------- EXERCISE 2 -------------
\section{}

Let $S$ be a perfect rigid body in motion with respect to a fixed reference frame $T$. $\vec{\Omega}$ is the instantaneous rotation vector of $S$ with respect to $T$.

Show that at a given instant, the projection of the velocities of the points of the solid on the axis $\vec{\Omega}$ is constant.

\begin{correctionbox}
Consider two arbitrary points $A$ and $B$ of the rigid body $S$. The velocity relationship for a rigid body is:
$$\vec{V}(B/T) = \vec{V}(A/T) + \vec{\Omega} \times \overrightarrow{AB}$$

Taking the projection on the axis of $\vec{\Omega}$ (scalar product with $\vec{\Omega}$):
$$\vec{V}(B/T) \cdot \vec{\Omega} = \vec{V}(A/T) \cdot \vec{\Omega} + (\vec{\Omega} \times \overrightarrow{AB}) \cdot \vec{\Omega}$$

The triple scalar product $(\vec{\Omega} \times \overrightarrow{AB}) \cdot \vec{\Omega}$ is zero because $\vec{\Omega} \times \overrightarrow{AB}$ is perpendicular to $\vec{\Omega}$:
$$(\vec{\Omega} \times \overrightarrow{AB}) \cdot \vec{\Omega} = 0$$

Therefore:
$$\vec{V}(B/T) \cdot \vec{\Omega} = \vec{V}(A/T) \cdot \vec{\Omega}$$

This shows that the projection of the velocity of any point $B$ on the axis $\vec{\Omega}$ equals the projection of the velocity of point $A$ on the same axis. Since $A$ and $B$ are arbitrary, this projection is constant for all points of the solid at a given instant.

\textbf{Physical interpretation:} All points of the solid have the same velocity component along the instantaneous rotation axis.
\end{correctionbox}

\newpage
% ----------- EXERCISE 3 -------------
\section{}

A rectangular plate $ABCD$ moves with respect to a direct orthonormal frame $R(O, X, Y, Z)$ such that $A$ remains coincident with point $O$ and side $AB$ remains in the plane $XOY$.

\begin{enumerate}
    \item Determine the number of degrees of freedom of the plate.
    \item What is the rotation vector $\vec{\Omega}$ of the plate with respect to $R$?
\end{enumerate}

\begin{correctionbox}
\begin{enumerate}
    \item \textbf{Number of degrees of freedom:}
    
    A rigid body in 3D space normally has 6 DOF (3 translations + 3 rotations).
    
    Constraints:
    \begin{itemize}
        \item Point $A$ is fixed at $O$: removes 3 DOF (translation blocked)
        \item Side $AB$ remains in plane $XOY$: This constrains the orientation. The plate can only rotate about point $A$ such that $AB$ stays in $XOY$. This removes 2 rotational DOF (no rotation about axes in the $XOY$ plane that would lift $AB$ out).
    \end{itemize}
    
    The plate can only rotate about an axis perpendicular to the $XOY$ plane, i.e., about the $Z$ axis.
    $$\text{DOF} = 6 - 3 - 2 = 1$$
    
    \item \textbf{Rotation vector:}
    
    Since the plate can only rotate about the $Z$ axis passing through $O$ (with $A$ at $O$ and $AB$ in $XOY$), the instantaneous rotation vector is:
    $$\vec{\Omega} = \dot{\theta} \vec{k}$$
    where $\theta$ is the angle of rotation about the $Z$ axis and $\vec{k}$ is the unit vector along $Z$.
    
    Alternatively, if we denote the angular velocity as $\omega$:
    $$\vec{\Omega} = \omega \vec{k}$$
\end{enumerate}
\end{correctionbox}

\newpage
% ----------- EXERCISE 4 -------------
\section{}

A hoop $S$ of center $C$ and radius $r$ moves in the plane $OX_1Y_1$ of a reference frame $R_1(O, X_1, Y_1, Z_1)$ while remaining in contact with the axis $OX_1$ at point $I$, $\overrightarrow{OI} = x\vec{i}_1$. The frame $R_1(O, X_1, Y_1, Z_1)$ is in motion with respect to the fixed reference frame $T_1(O_1, X_1, Y_1, Z_1)$ such that $\overrightarrow{O_1O} = y\vec{j}_1$.

Let $M$ be a point of the solid $S$. We set $(\vec{i}_1, \overrightarrow{CM}) = \theta$. We attach to the solid a frame $T(C, X, Y, Z)$ with basis $(\vec{i}, \vec{j}, \vec{k})$ such that the $X$ axis passes through point $M$.

\begin{enumerate}
    \item Determine the number of degrees of freedom of the solid.
    \item Express the instantaneous rotation vector of the hoop with respect to $T_1$.
    \item Can there be pivoting of the hoop with respect to $R$?
    \item Determine the sliding velocity of the hoop with respect to the axis $OX_1$.
    \item Calculate $\vec{V}(I/R)$ and $\vec{V}(I/S)$ then find again the sliding velocity of the hoop with respect to the axis $OX_1$.
    \item Calculate the acceleration $\vec{\gamma}(M/T_1)$. Deduce the accelerations $\vec{\gamma}_S(I/T_1)$ and $\vec{\gamma}_S(J/T_1)$ where $J$ is a point diametrically opposite to $I$ on $S$.
\end{enumerate}

\begin{correctionbox}
\begin{enumerate}
    \item \textbf{Degrees of freedom:}
    
    The hoop moves in plane $OX_1Y_1$ and remains in contact with axis $OX_1$. 
    \begin{itemize}
        \item Position of center $C$: $x$ (along $X_1$) and fixed at $y_C = r$ (above $X_1$), so 1 DOF
        \item Rotation in plane: angle $\theta$, so 1 DOF
    \end{itemize}
    Also, frame $R_1$ moves relative to $T_1$: position $y$ of $O$ relative to $O_1$, so 1 DOF.
    $$\text{Total DOF} = 3 \text{ (positions: } x, y, \theta)$$
    
    \item \textbf{Rotation vector of hoop with respect to $T_1$:}
    
    Using the composition of angular velocities:
    $$\vec{\Omega}(S/T_1) = \vec{\Omega}(S/R_1) + \vec{\Omega}(R_1/T_1)$$
    
    Since the hoop rotates in the $X_1Y_1$ plane:
    $$\vec{\Omega}(S/R_1) = \dot{\theta} \vec{k}_1$$
    
    Frame $R_1$ only translates relative to $T_1$ (no rotation):
    $$\vec{\Omega}(R_1/T_1) = \vec{0}$$
    
    Therefore:
    $$\vec{\Omega}(S/T_1) = \dot{\theta} \vec{k}_1$$
    
    \item \textbf{Pivoting with respect to $R_1$:}
    
    Pivoting means rotation about an axis normal to the plane of motion. Since the hoop remains in plane $OX_1Y_1$ and $\vec{\Omega}(S/R_1) = \dot{\theta} \vec{k}_1$ is perpendicular to this plane, there is pivoting about the $Z_1$ axis.
    
    \textbf{Yes, there is pivoting.}
    
    \item \textbf{Sliding velocity with respect to axis $OX_1$:}
    
    The sliding velocity is the velocity of the contact point $I$ along the axis $OX_1$ in frame $R_1$:
    $$v_{\text{slide}} = \dot{x}$$
    
    Since $\overrightarrow{OI} = x\vec{i}_1$ and $O$ is fixed in $R_1$:
    $$\vec{V}(I/R_1) = \dot{x} \vec{i}_1 + \vec{\Omega}(S/R_1) \times \overrightarrow{CI}$$
    
    With $\overrightarrow{CI} = -r\vec{j}_1$ (from center to contact point):
    $$\vec{\Omega}(S/R_1) \times \overrightarrow{CI} = \dot{\theta}\vec{k}_1 \times (-r\vec{j}_1) = r\dot{\theta}\vec{i}_1$$
    
    For rolling without slipping: $\vec{V}(I/R_1) = \vec{0}$, giving $\dot{x} + r\dot{\theta} = 0$, so $\dot{x} = -r\dot{\theta}$.
    
    The sliding velocity is:
    $$v_{\text{slide}} = \dot{x}$$
    
    \item \textbf{Velocities $\vec{V}(I/R_1)$ and $\vec{V}(I/S)$:}
    
    Point $I$ as part of the solid $S$:
    $$\vec{V}(I/S) = \vec{0}$$ (by definition, $I$ is at rest in $S$)
    
    Actually, $I$ is the instantaneous contact point, so:
    $$\vec{V}(I/S) = \vec{V}(C/R_1) + \vec{\Omega}(S/R_1) \times \overrightarrow{CI}$$
    
    With $\vec{V}(C/R_1) = \dot{x}\vec{i}_1$ and $\overrightarrow{CI} = -r\vec{j}_1$:
    $$\vec{V}(I/S) = \dot{x}\vec{i}_1 + \dot{\theta}\vec{k}_1 \times (-r\vec{j}_1) = \dot{x}\vec{i}_1 + r\dot{\theta}\vec{i}_1 = (\dot{x} + r\dot{\theta})\vec{i}_1$$
    
    The sliding velocity is the component of $\vec{V}(I/R_1)$ along $\vec{i}_1$:
    $$v_{\text{slide}} = \dot{x} + r\dot{\theta}$$
    
    \item \textbf{Acceleration $\vec{\gamma}(M/T_1)$:}
    
    Using the general formula:
    $$\vec{\gamma}(M/T_1) = \vec{\gamma}(C/T_1) + \frac{d\vec{\Omega}}{dt}(S/T_1) \times \overrightarrow{CM} + \vec{\Omega}(S/T_1) \times (\vec{\Omega}(S/T_1) \times \overrightarrow{CM})$$
    
    With $\vec{\gamma}(C/T_1) = \ddot{x}\vec{i}_1 + \ddot{y}\vec{j}_1$, $\vec{\Omega}(S/T_1) = \dot{\theta}\vec{k}_1$, and $\overrightarrow{CM} = r\vec{i}$ (where $\vec{i}$ is the radial direction to $M$):
    
    In the $R_1$ basis: $\overrightarrow{CM} = r\cos\theta\vec{i}_1 + r\sin\theta\vec{j}_1$
    
    The centripetal acceleration term:
    $$\vec{\Omega} \times (\vec{\Omega} \times \overrightarrow{CM}) = -\dot{\theta}^2 r(\cos\theta\vec{i}_1 + \sin\theta\vec{j}_1)$$
    
    Therefore:
    $$\vec{\gamma}(M/T_1) = (\ddot{x} - r\dot{\theta}^2\cos\theta)\vec{i}_1 + (\ddot{y} - r\dot{\theta}^2\sin\theta)\vec{j}_1 + r\ddot{\theta}(-\sin\theta\vec{i}_1 + \cos\theta\vec{j}_1)$$
    
    For point $I$ (at $\theta = -\pi/2$, i.e., bottom of hoop):
    $$\vec{\gamma}_S(I/T_1) = (\ddot{x} + r\ddot{\theta})\vec{i}_1 + (\ddot{y} + r\dot{\theta}^2)\vec{j}_1$$
    
    For point $J$ (at $\theta = \pi/2$, top of hoop):
    $$\vec{\gamma}_S(J/T_1) = (\ddot{x} - r\ddot{\theta})\vec{i}_1 + (\ddot{y} - r\dot{\theta}^2)\vec{j}_1$$
\end{enumerate}
\end{correctionbox}

\newpage
% ----------- EXERCISE 5 -------------
\section{}

Consider a homogeneous sphere $S(G, a)$, of center $G$, mass $m$, and radius $a$, moving on the fixed horizontal plane $(O, X_0, Y_0)$ of a direct orthonormal fixed frame $R_0(O_0, X_0, Y_0, Z_0)$. Let $R_S(G, X, Y, Z)$ be the direct orthonormal frame attached to the sphere. $I$ is the contact point between the sphere $S$ and the plane $(O, X_0, Y_0)$. We denote by $\psi$, $\theta$, and $\varphi$ the three Euler angles of $S/R_0$ and $x$, $y$, and $z$ the coordinates of $G$. We assume that $S$ rolls without slipping on the axis $(O, X_0)$.

\textbf{1. Euler angles}
\begin{enumerate}
    \item[a)] Identify the relative Euler frames $R_1$ and $R_2$. Deduce the angular variables.
    \item[b)] Parameterize the sphere $S$.
    \item[c)] Determine the expression of the instantaneous rotation vector of $S$ with respect to $R_0$.
\end{enumerate}

\textbf{2. Kinematics of the contact point}
\begin{enumerate}
    \item[a)] Express the non-pivoting condition of $S$ with respect to $R_0$.
    \item[b)] Express the non-rolling condition of $S$ with respect to $(O, X_0)$. Show that one of the obtained equations is semi-holonomic.
    \item[c)] Calculate the velocities of the different contact points: Point of the sphere $I \in S$; point of the material plane $I \in (O, X_0, Y_0)$; and geometric point $I \in E$.
    \item[d)] Give the no-slip condition of $S$ with respect to $(O, X_0, Y_0)$. What is the number of degrees of freedom of $S$ in its motion with respect to $R_0$?
\end{enumerate}

\begin{correctionbox}
\textbf{1. Euler angles}

\begin{enumerate}
    \item[a)] \textbf{Euler frames:}
    
    The Euler angles $(\psi, \theta, \varphi)$ define the orientation:
    \begin{itemize}
        \item $R_1$: obtained from $R_0$ by rotation $\psi$ about $Z_0$ (precession)
        \item $R_2$: obtained from $R_1$ by rotation $\theta$ about the line of nodes (nutation)
        \item $R_S$: obtained from $R_2$ by rotation $\varphi$ about $Z$ (spin)
    \end{itemize}
    
    Angular variables: $\psi$ (precession), $\theta$ (nutation), $\varphi$ (spin)
    
    \item[b)] \textbf{Parameterization:}
    
    The sphere requires 6 parameters:
    \begin{itemize}
        \item Position of center $G$: $(x, y, z)$ in $R_0$
        \item Orientation: Euler angles $(\psi, \theta, \varphi)$
    \end{itemize}
    
    Since the sphere rolls on plane $(O, X_0, Y_0)$: $z = a$ (constant).
    
    Parameters: $(x, y, \psi, \theta, \varphi)$ with constraint $z = a$.
    
    \item[c)] \textbf{Instantaneous rotation vector:}
    
    Using Euler angles:
    $$\vec{\Omega}(S/R_0) = \dot{\psi}\vec{k}_0 + \dot{\theta}\vec{u} + \dot{\varphi}\vec{k}$$
    where $\vec{u}$ is the line of nodes and $\vec{k}$ is the axis of $R_S$.
    
    In the $R_0$ frame:
    $$\vec{\Omega}(S/R_0) = \dot{\psi}\vec{k}_0 + \dot{\theta}\vec{u} + \dot{\varphi}\vec{k}$$
\end{enumerate}

\textbf{2. Kinematics of contact point}

\begin{enumerate}
    \item[a)] \textbf{Non-pivoting condition:}
    
    Non-pivoting means no rotation about the vertical axis at the contact point. The vertical component of angular velocity must be zero:
    $$\vec{\Omega}(S/R_0) \cdot \vec{k}_0 = 0$$
    
    This gives: $\dot{\psi} + \dot{\varphi}\cos\theta = 0$
    
    \item[b)] \textbf{Non-rolling condition on $(O, X_0)$:}
    
    Rolling on axis $X_0$ means the sphere moves along $X_0$ without slipping perpendicular to it. The velocity of contact point $I$ must be along $X_0$:
    $$\vec{V}(I/R_0) \cdot \vec{j}_0 = 0$$
    
    With $\vec{V}(I/R_0) = \vec{V}(G/R_0) + \vec{\Omega}(S/R_0) \times \overrightarrow{GI}$:
    $$\dot{y} + [\vec{\Omega}(S/R_0) \times (-a\vec{k}_0)]_y = 0$$
    
    This yields: $\dot{y} - a\omega_x = 0$ where $\omega_x$ is the $X_0$ component of $\vec{\Omega}$.
    
    One equation is \textbf{semi-holonomic} if it involves velocities but cannot be integrated to give a position constraint (e.g., $\dot{y} = a\dot{\varphi}\sin\psi$).
    
    \item[c)] \textbf{Velocities:}
    
    \begin{itemize}
        \item Point $I \in S$ (material point of sphere):
        $$\vec{V}(I \in S/R_0) = \vec{V}(G/R_0) + \vec{\Omega}(S/R_0) \times \overrightarrow{GI}$$
        $$= \dot{x}\vec{i}_0 + \dot{y}\vec{j}_0 + \vec{\Omega} \times (-a\vec{k}_0)$$
        
        \item Point $I \in$ plane (material point of fixed plane):
        $$\vec{V}(I \in \text{plane}/R_0) = \vec{0}$$
        
        \item Geometric point $I \in E$ (space point):
        $$\vec{V}(I \in E/R_0) = \vec{V}(G/R_0) = \dot{x}\vec{i}_0 + \dot{y}\vec{j}_0$$
    \end{itemize}
    
    \item[d)] \textbf{No-slip condition and DOF:}
    
    No-slip means $\vec{V}(I \in S/R_0) = \vec{V}(I \in \text{plane}/R_0) = \vec{0}$:
    $$\dot{x}\vec{i}_0 + \dot{y}\vec{j}_0 + \vec{\Omega}(S/R_0) \times (-a\vec{k}_0) = \vec{0}$$
    
    This gives 2 constraints (in $X_0$ and $Y_0$ directions). Combined with rolling on $X_0$ (1 additional constraint):
    
    Total parameters: 5 (since $z = a$). Constraints: 2 (no-slip).
    $$\text{DOF} = 5 - 2 = 3$$
\end{enumerate}
\end{correctionbox}

\newpage
% ----------- EXERCISE 6 -------------
\section{}

A horizontal disk $S_2$, of center $O_1$ and radius $R$, is rotated about the axis $O_1Z_1$ of the reference frame $T_1(O_1, X_1, Y_1, Z_1)$ assumed fixed, with angular velocity $\omega$, by a disk $S_1$, of center $C$ and radius $r$. $S_1$ rolls without slipping on the contour of disk $S_2$, with angular velocity $\Omega$, about the rod $AC$ which is parallel to the plane of $S_2$. This rod rotates about the axis $AZ_1$ with angular velocity $\omega_0$ with respect to $S_2$.

\begin{enumerate}
    \item Write the expressions for the rotation velocity vectors $\vec{\Omega}(S_1/AC)$, $\vec{\Omega}(AC/S_2)$, and $\vec{\Omega}(S_2/T_1)$.
    \item By exploiting the rolling without slipping condition of $S_1$ on $S_2$, deduce a relation between the angular velocities $\Omega$, $\omega_0$, and $\omega$.
    \item Express the velocity of point $J$ of $S_1$, which is diametrically opposite to $I$. What does this velocity become if disk $S_2$ is stationary?
\end{enumerate}

\begin{correctionbox}
\begin{enumerate}
    \item \textbf{Rotation velocity vectors:}
    
    \begin{itemize}
        \item $\vec{\Omega}(S_1/AC)$: Disk $S_1$ rotates about rod $AC$ with angular velocity $\Omega$. If $\vec{u}_{AC}$ is the unit vector along $AC$:
        $$\vec{\Omega}(S_1/AC) = \Omega \vec{u}_{AC}$$
        
        \item $\vec{\Omega}(AC/S_2)$: Rod $AC$ rotates about axis $AZ_1$ with angular velocity $\omega_0$ relative to $S_2$:
        $$\vec{\Omega}(AC/S_2) = \omega_0 \vec{k}_1$$
        
        \item $\vec{\Omega}(S_2/T_1)$: Disk $S_2$ rotates about $O_1Z_1$ with angular velocity $\omega$:
        $$\vec{\Omega}(S_2/T_1) = \omega \vec{k}_1$$
    \end{itemize}
    
    \item \textbf{Relation between angular velocities:}
    
    Using composition of angular velocities:
    $\vec{\Omega}(S_1/T_1) = \vec{\Omega}(S_1/AC) + \vec{\Omega}(AC/S_2) + \vec{\Omega}(S_2/T_1)$
    $= \Omega \vec{u}_{AC} + \omega_0 \vec{k}_1 + \omega \vec{k}_1 = \Omega \vec{u}_{AC} + (\omega_0 + \omega) \vec{k}_1$
    
    At the contact point $I$ between $S_1$ and $S_2$, the rolling without slipping condition requires:
    $\vec{V}(I \in S_1/T_1) = \vec{V}(I \in S_2/T_1)$
    
    For disk $S_2$, point $I$ is at radius $R$ from $O_1$:
    $\vec{V}(I \in S_2/T_1) = \vec{\Omega}(S_2/T_1) \times \overrightarrow{O_1I} = \omega \vec{k}_1 \times R\vec{u}_r$
    where $\vec{u}_r$ is the radial direction from $O_1$ to $I$.
    
    For disk $S_1$, point $I$ is at the contact (on the rim):
    $\vec{V}(I \in S_1/T_1) = \vec{V}(C/T_1) + \vec{\Omega}(S_1/T_1) \times \overrightarrow{CI}$
    
    The center $C$ of $S_1$ moves in a circle of radius $R$ around $O_1$:
    $\vec{V}(C/T_1) = (\omega_0 + \omega) \vec{k}_1 \times R\vec{u}_r$
    
    At the contact point, $\overrightarrow{CI}$ is perpendicular to $AC$ with magnitude $r$. The no-slip condition gives:
    $\vec{V}(I \in S_1/T_1) = \vec{V}(I \in S_2/T_1)$
    
    Expanding and comparing tangential components:
    $(\omega_0 + \omega)R - \Omega r = \omega R$
    
    Simplifying:
    $\omega_0 R = \Omega r$
    
    Therefore:
    $\boxed{\Omega = \frac{\omega_0 R}{r}}$
    
    Or including the motion of $S_2$:
    $\boxed{\Omega = \frac{(\omega_0 + \omega)R - \omega R}{r} = \frac{\omega_0 R}{r}}$
    
    \item \textbf{Velocity of point $J$:}
    
    Point $J$ is diametrically opposite to $I$ on disk $S_1$. Thus $\overrightarrow{CJ} = -\overrightarrow{CI}$.
    
    The velocity of $J$ is:
    $\vec{V}(J/T_1) = \vec{V}(C/T_1) + \vec{\Omega}(S_1/T_1) \times \overrightarrow{CJ}$
    $= (\omega_0 + \omega) \vec{k}_1 \times R\vec{u}_r + [\Omega \vec{u}_{AC} + (\omega_0 + \omega) \vec{k}_1] \times (-\overrightarrow{CI})$
    
    Since $\vec{\Omega}(S_1/T_1) \times \overrightarrow{CI}$ gives the velocity contribution at $I$, and we have:
    $\vec{\Omega}(S_1/T_1) \times \overrightarrow{CJ} = -\vec{\Omega}(S_1/T_1) \times \overrightarrow{CI}$
    
    From the contact point analysis:
    $\vec{V}(I/T_1) = \vec{V}(C/T_1) + \vec{\Omega}(S_1/T_1) \times \overrightarrow{CI} = \omega R \vec{u}_{\theta}$
    
    Therefore:
    $\vec{V}(J/T_1) = \vec{V}(C/T_1) - [\vec{V}(I/T_1) - \vec{V}(C/T_1)]$
    $= 2\vec{V}(C/T_1) - \vec{V}(I/T_1)$
    $= 2(\omega_0 + \omega)R\vec{u}_{\theta} - \omega R\vec{u}_{\theta}$
    $= [(2\omega_0 + 2\omega - \omega)R]\vec{u}_{\theta}$
    $\boxed{\vec{V}(J/T_1) = (2\omega_0 + \omega)R\vec{u}_{\theta}}$
    
    \textbf{If disk $S_2$ is stationary} ($\omega = 0$):
    $\boxed{\vec{V}(J/T_1) = 2\omega_0 R\vec{u}_{\theta}}$
    
    The point $J$ moves with twice the velocity of the center $C$, which is characteristic of rolling motion.
\end{enumerate}
\end{correctionbox}

\end{document}